\chapter{Tổng hợp, đánh giá và thảo luận kết quả kiểm nghiệm}
Chương này tổng hợp số liệu đã ghi nhận trong quá trình thực tập, đánh giá độ lặp lại của từng phép đo và tạo cơ sở để đối chiếu với chỉ tiêu chất lượng công bố của từng sản phẩm.

\section{Kết quả phân tích và đánh giá độ lặp lại}
Độ lặp lại được xem xét trước hết thông qua chênh lệch tuyệt đối $d$ và chênh lệch tương đối giữa hai kết quả (RPD):
\[
  d=\left|x_1-x_2\right|, \qquad
  \mathrm{RPD}(\%)=\frac{\left|x_1-x_2\right|}{\left(x_1+x_2\right)/2}\times100.
\]
Giá trị RPD của từng mẫu được đối chiếu với các giới hạn chấp nhận trong quy định kiểm soát nội bộ của công ty, được tổng hợp tại Bảng~\ref{tab:nguong-rpd}.

\begin{table}[H]
  \centering
  \small
  \setlength{\tabcolsep}{4pt}
  \renewcommand{\arraystretch}{1.15}
  \caption{Ngưỡng chấp nhận độ lặp lại cho các chỉ tiêu kiểm nghiệm}
  \label{tab:nguong-rpd}
  \begin{tabular}{|>{\centering\arraybackslash}p{0.08\textwidth}|p{0.25\textwidth}|>{\centering\arraybackslash}p{0.28\textwidth}|>{\centering\arraybackslash}p{0.29\textwidth}|}
    \hline
    \textbf{STT} & \textbf{Chỉ tiêu kiểm nghiệm} & \textbf{Phương pháp/TCVN tham khảo} & \textbf{Tiêu chí lặp lại áp dụng} \\
    \hline
    1 & Đạm tổng số & TCVN~8557:2010 & $\mathrm{RPD}\leq5\%$ \\
    \hline
    2 & Lân hữu hiệu (\(\mathrm{P_2O_5}\)) & TCVN~8559:2010 & $\mathrm{RPD}\leq5\%$ \\
    \hline
    3 & Kali hữu hiệu (\(\mathrm{K_2O}\)) & TCVN~8560:2018 & $\mathrm{RPD}\leq5\%$ theo SOP nội bộ \\
    \hline
    4 & Carbon hữu cơ tổng số & TCVN~9294:2012, Walkley--Black & $\mathrm{RPD}\leq10\%$ \\
    \hline
    5 & Độ ẩm (\(\mathrm{H_2O}\)) & TCVN~9297:2012 & Không đánh giá bằng RPD; $d\leq0{,}3$ điểm \% \\
    \hline
  \end{tabular}
\end{table}


\noindent\textit{Nguồn: Tài liệu nội bộ của Công ty TNHH Quốc Tế Khánh Sinh \cite{company-internal-document}.}

\subsection{Kết quả xác định đạm tổng số}
\NitrogenResultsTable

\subsubsection{Đánh giá độ lặp lại}
Kết quả ở Bảng~\ref{tab:ket-qua-nito} được tổng hợp lại theo từng ngày và từng mã mẫu trong Bảng~\ref{tab:do-lap-lai-nito}. Sai lệch tương đối giữa hai lần phân tích được tính theo RPD; tiêu chí chấp nhận áp dụng cho phép đo đạm tổng số là $\mathrm{RPD}\leq5\%$.

\clearpage
{\scriptsize
\setlength{\tabcolsep}{3pt}
\renewcommand{\arraystretch}{1.2}
\begin{longtable}{|>{\centering\arraybackslash}p{0.14\textwidth}|>{\centering\arraybackslash}p{0.12\textwidth}|>{\centering\arraybackslash}p{0.13\textwidth}|>{\centering\arraybackslash}p{0.13\textwidth}|>{\centering\arraybackslash}p{0.13\textwidth}|>{\centering\arraybackslash}p{0.13\textwidth}|>{\centering\arraybackslash}p{0.10\textwidth}|}
\caption{Đánh giá độ lặp lại phép xác định đạm tổng số theo từng mẫu}\label{tab:do-lap-lai-nito} \\
\hline
\textbf{Ngày} & \textbf{KHM} & \textbf{Lần 1 (\%N)} & \textbf{Lần 2 (\%N)} & \textbf{$d$ (điểm \%)} & \textbf{RPD (\%)} & \textbf{Đánh giá} \\
\hline
\endfirsthead
\multicolumn{7}{c}{\tablename\ \thetable{} -- tiếp theo} \\
\hline
\textbf{Ngày} & \textbf{KHM} & \textbf{Lần 1 (\%N)} & \textbf{Lần 2 (\%N)} & \textbf{$d$ (điểm \%)} & \textbf{RPD (\%)} & \textbf{Đánh giá} \\
\hline
\endhead
\hline
\multicolumn{7}{r}{\textit{Còn tiếp trang sau}} \\
\endfoot
\hline
\endlastfoot
01/07/2026 & 501X1 & 12,16 & 12,46 & 0,30 & 2,44 & Đạt \\
\hline
03/07/2026 & 603X1 & 14,36 & 14,22 & 0,14 & 0,98 & Đạt \\
\hline
03/07/2026 & 401X2 & 5,88 & 5,88 & 0,00 & 0,00 & Đạt \\
\hline
03/07/2026 & 400X1 & 19,87 & 19,87 & 0,00 & 0,00 & Đạt \\
\hline
03/07/2026 & 300X1 & 9,56 & 9,56 & 0,00 & 0,00 & Đạt \\
\hline
03/07/2026 & 403X1 & 14,82 & 14,95 & 0,13 & 0,87 & Đạt \\
\hline
04/07/2026 & 401X1 & 5,64 & 5,69 & 0,05 & 0,88 & Đạt \\
\hline
04/07/2026 & 20X3 & 19,62 & 19,76 & 0,14 & 0,71 & Đạt \\
\hline
05/07/2026 & 507X1 & 13,13 & 13,38 & 0,25 & 1,89 & Đạt \\
\hline
05/07/2026 & 02X1(tr) & 20,06 & 20,18 & 0,12 & 0,60 & Đạt \\
\hline
06/07/2026 & 02X1(v) & 20,32 & 20,45 & 0,13 & 0,64 & Đạt \\
\hline
06/07/2026 & 01X1 & 24,39 & 24,46 & 0,07 & 0,29 & Đạt \\
\hline
06/07/2026 & 12X1 & 11,78 & 11,91 & 0,13 & 1,10 & Đạt \\
\hline
06/07/2026 & DAP & 14,72 & 14,85 & 0,13 & 0,88 & Đạt \\
\hline
08/07/2026 & 05X1 & 4,94 & 5,04 & 0,10 & 2,00 & Đạt \\
\hline
11/07/2026 & 20X3 & 20,19 & 20,05 & 0,14 & 0,70 & Đạt \\
\hline
11/07/2026 & 507X1 & 13,46 & 13,59 & 0,13 & 0,96 & Đạt \\
\hline
11/07/2026 & 01X1 & 24,45 & 24,68 & 0,23 & 0,94 & Đạt \\
\hline
13/07/2026 & 46X1 & 44,76 & 44,90 & 0,14 & 0,31 & Đạt \\
\hline
14/07/2026 & 400X1 & 20,34 & 20,47 & 0,13 & 0,64 & Đạt \\
\hline
14/07/2026 & 403X1 & 15,57 & 15,57 & 0,00 & 0,00 & Đạt \\
\hline
14/07/2026 & 300X1 & 9,22 & 9,35 & 0,13 & 1,40 & Đạt \\
\hline
14/07/2026 & 309X1 & 14,37 & 14,50 & 0,13 & 0,90 & Đạt \\
\hline
15/07/2026 & 507X1 & 14,10 & 14,24 & 0,14 & 0,99 & Đạt \\
\hline
16/07/2026 & 403X1 & 14,39 & 14,52 & 0,13 & 0,90 & Đạt \\
\hline
16/07/2026 & 02X1 & 20,23 & 20,37 & 0,14 & 0,69 & Đạt \\
\hline
16/07/2026 & 46X1 & 45,97 & 45,97 & 0,00 & 0,00 & Đạt \\
\hline
18/07/2026 & 02X1 & 20,06 & 20,06 & 0,00 & 0,00 & Đạt \\
\hline
20/07/2026 & 02X1 & 17,78 & 18,10 & 0,32 & 1,78 & Đạt \\
\hline
21/07/2026 & 402X1 & 12,49 & 12,61 & 0,12 & 0,96 & Đạt \\
\hline
25/07/2026 & 20X3 & 18,05 & 17,92 & 0,13 & 0,72 & Đạt \\
\hline
25/07/2026 & 02X1 & 20,28 & 20,40 & 0,12 & 0,59 & Đạt \\
\hline
26/07/2026 & 20X3 & 18,85 & 18,98 & 0,13 & 0,69 & Đạt \\
\hline
26/07/2026 & 02X1 & 20,43 & 20,43 & 0,00 & 0,00 & Đạt \\
\hline
27/07/2026 & 02X1 & 20,38 & 20,25 & 0,13 & 0,64 & Đạt \\
\hline
27/07/2026 & 01X1 & 24,50 & 24,64 & 0,14 & 0,57 & Đạt \\
\hline
27/07/2026 & 20X3 & 19,10 & 19,03 & 0,07 & 0,37 & Đạt \\
\hline
27/07/2026 & 408X3 & 9,95 & 10,09 & 0,14 & 1,40 & Đạt \\
\hline
28/07/2026 & 01X1 & 24,88 & 25,17 & 0,29 & 1,16 & Đạt \\
\hline
29/07/2026 & 01X1 & 25,06 & 25,06 & 0,00 & 0,00 & Đạt \\
\hline
30/07/2026 & 01X1 & 24,83 & 24,83 & 0,00 & 0,00 & Đạt \\
\hline
\end{longtable}

}

\noindent\textbf{Nhận xét:} Kết quả đánh giá độ lặp của phép xác định hàm lượng nitơ cho thấy tất cả 14 mẫu đều có giá trị RPD nhỏ hơn 5\%, đáp ứng tiêu chí chấp nhận về độ lặp của phương pháp.
Giá trị RPD dao động từ 0,00\% đến 2,44\%, trong đó giá trị lớn nhất ghi nhận ở mẫu 501X1 ngày 01/06/2026 là 2,44\%, vẫn thấp hơn đáng kể so với giới hạn cho phép.
Giá trị RPD trung bình của tập mẫu vào khoảng 0,81\%, cho thấy mức độ sai khác giữa hai lần phân tích nhìn chung rất nhỏ.
Nhiều mẫu có RPD bằng 0 hoặc dưới 1\%, chứng tỏ kết quả giữa các lần lặp có độ tương đồng cao. 
Tuy nhiên, kết quả này chủ yếu phản ánh độ lặp trong cùng điều kiện phân tích vì vậy chưa đủ để kết luận về độ đúng hay độ tái lập của phương pháp.

\subsection{Kết quả xác định kali hữu hiệu}
\PotassiumResultsTable

\subsubsection{Đánh giá độ lặp lại}
Các kết quả \%K$_2$O ở Bảng~\ref{tab:ket-qua-kali} được tổng hợp theo từng ngày và mã mẫu trong Bảng~\ref{tab:do-lap-lai-kali}. Tiêu chí chấp nhận đối với phép xác định kali hữu hiệu là $\mathrm{RPD}\leq5\%$.

{\scriptsize
\setlength{\tabcolsep}{3pt}
\renewcommand{\arraystretch}{1.2}
\begin{longtable}{|>{\centering\arraybackslash}p{0.14\textwidth}|>{\centering\arraybackslash}p{0.12\textwidth}|>{\centering\arraybackslash}p{0.13\textwidth}|>{\centering\arraybackslash}p{0.13\textwidth}|>{\centering\arraybackslash}p{0.13\textwidth}|>{\centering\arraybackslash}p{0.13\textwidth}|>{\centering\arraybackslash}p{0.10\textwidth}|}
\caption{Đánh giá độ lặp lại phép xác định kali hữu hiệu theo từng mẫu}\label{tab:do-lap-lai-kali} \\
\hline
\textbf{Ngày} & \textbf{KHM} & \textbf{Lần 1 (\%K$_2$O)} & \textbf{Lần 2 (\%K$_2$O)} & \textbf{$d$ (điểm \%)} & \textbf{RPD (\%)} & \textbf{Đánh giá} \\
\hline
\endfirsthead
\multicolumn{7}{c}{\tablename\ \thetable{} -- tiếp theo} \\
\hline
\textbf{Ngày} & \textbf{KHM} & \textbf{Lần 1 (\%K$_2$O)} & \textbf{Lần 2 (\%K$_2$O)} & \textbf{$d$ (điểm \%)} & \textbf{RPD (\%)} & \textbf{Đánh giá} \\
\hline
\endhead
\hline
\multicolumn{7}{r}{\textit{Còn tiếp trang sau}} \\
\endfoot
\hline
\endlastfoot
01/07/2026 & 501X1 & 13,39 & 13,24 & 0,15 & 1,13 & Đạt \\
\hline
03/07/2026 & 603X1 & 8,65 & 9,12 & 0,47 & 5,29 & Không đạt \\
\hline
03/07/2026 & 401X2 & 3,20 & 3,20 & 0,00 & 0,00 & Đạt \\
\hline
03/07/2026 & 400X1 & 12,78 & 12,18 & 0,60 & 4,81 & Đạt \\
\hline
03/07/2026 & 300X1 & 3,14 & 3,14 & 0,00 & 0,00 & Đạt \\
\hline
03/07/2026 & 403X1 & 5,58 & 4,39 & 1,19 & 23,87 & Không đạt \\
\hline
04/07/2026 & 401X1 & 3,25 & 3,25 & 0,00 & 0,00 & Đạt \\
\hline
04/07/2026 & 20X3 & 10,75 & 10,56 & 0,19 & 1,78 & Đạt \\
\hline
05/07/2026 & 507X1 & 7,39 & 7,46 & 0,07 & 0,94 & Đạt \\
\hline
06/07/2026 & 12X1 & 9,81 & 9,29 & 0,52 & 5,45 & Không đạt \\
\hline
08/07/2026 & 05X1 & 2,86 & 2,86 & 0,00 & 0,00 & Đạt \\
\hline
11/07/2026 & 20X3 & 9,75 & 9,88 & 0,13 & 1,32 & Đạt \\
\hline
11/07/2026 & 507X1 & 7,44 & 7,45 & 0,01 & 0,13 & Đạt \\
\hline
16/07/2026 & 403X1 & 5,54 & 5,67 & 0,13 & 2,32 & Đạt \\
\hline
16/07/2026 & 94MX1 & 60,44 & 62,99 & 2,55 & 4,13 & Đạt \\
\hline
21/07/2026 & 402X1 & 7,63 & 7,21 & 0,42 & 5,66 & Không đạt \\
\hline
25/07/2026 & 20X3 & 9,64 & 9,60 & 0,04 & 0,42 & Đạt \\
\hline
26/07/2026 & 20X3 & 10,31 & 10,56 & 0,25 & 2,40 & Đạt \\
\hline
27/07/2026 & 408X3 & 25,21 & 25,07 & 0,14 & 0,56 & Đạt \\
\hline
27/07/2026 & 04X1 & 61,95 & 60,19 & 1,76 & 2,88 & Đạt \\
\hline
\end{longtable}

}

\noindent\textbf{Nhận xét:} Kết quả đánh giá độ lặp của phép xác định hàm lượng kali hữu hiệu cho thấy giá trị RPD dao động từ 0,00\% đến 8,66\%.
Với tiêu chí chấp nhận RPD $\leq$ 5\%, có 8/10 mẫu đạt yêu cầu, chiếm 80,0\%, và 2/10 mẫu không đạt, chiếm 20,0\%.
Giá trị RPD trung bình của tập mẫu khoảng 2,80\%, cho thấy mức sai khác giữa hai lần phân tích nhìn chung tương đối nhỏ.

Hai mẫu không đáp ứng yêu cầu gồm 403X1 ngày 03/06/2026 (RPD = 8,66\%) và 12X1 ngày 06/06/2026 (RPD = 5,18\%).
Trong đó, mẫu 403X1 có mức sai khác lớn nhất, với kết quả hai lần đo lần lượt là 5,78\% và 5,30\% K$_2$O, cần được ưu tiên kiểm tra và phân tích lại. Mẫu 12X1 chỉ vượt nhẹ giới hạn 5\%, tuy nhiên vẫn được đánh giá là không đạt theo tiêu chí đã đặt ra.

Đáng chú ý, một số mẫu như 603X1 (RPD = 4,99\%) và 400X1 (RPD = 4,63\%) tuy vẫn đạt yêu cầu nhưng có giá trị khá gần giới hạn cho phép, do đó nên được theo dõi trong các lần phân tích tiếp theo.
Ngược lại, nhiều mẫu có RPD rất thấp cho thấy khả năng lặp lại tốt trong điều kiện phân tích tương ứng.

Nhìn chung, phép xác định kali hữu hiệu bằng phương pháp quang kế ngọn lửa cho kết quả có độ lặp tương đối tốt, nhưng vẫn xuất hiện một số trường hợp sai khác vượt giới hạn.
Đối với các mẫu không đạt, cần tiến hành đồng nhất và chuẩn bị lại mẫu, thực hiện lại quá trình chiết và pha loãng, đồng thời kiểm tra đường chuẩn, độ ổn định của ngọn lửa và hệ thống hút mẫu của máy quang kế ngọn lửa trước khi đưa ra kết quả cuối cùng.

\subsection{Kết quả xác định lân hữu hiệu}
\PhosphorusResultsTable

\subsubsection{Đánh giá độ lặp lại}
Sự phù hợp giữa hai lần lặp của phép xác định lân hữu hiệu được đánh giá từ kết quả \%P$_2$O$_5$ cuối cùng. Bảng~\ref{tab:do-lap-lai-lan} sử dụng tiêu chí $\mathrm{RPD}\leq5\%$.

{\scriptsize
\setlength{\tabcolsep}{3pt}
\renewcommand{\arraystretch}{1.2}
\begin{longtable}{|>{\centering\arraybackslash}p{0.14\textwidth}|>{\centering\arraybackslash}p{0.12\textwidth}|>{\centering\arraybackslash}p{0.13\textwidth}|>{\centering\arraybackslash}p{0.13\textwidth}|>{\centering\arraybackslash}p{0.13\textwidth}|>{\centering\arraybackslash}p{0.13\textwidth}|>{\centering\arraybackslash}p{0.10\textwidth}|}
\caption{Đánh giá độ lặp lại phép xác định lân hữu hiệu trong tuần 01--06/06/2026}\label{tab:do-lap-lai-lan} \\
\hline
\textbf{Ngày} & \textbf{KHM} & \textbf{Lần 1 (\(\%\mathrm{P_2O_5}\))} & \textbf{Lần 2 (\(\%\mathrm{P_2O_5}\))} & \textbf{$d$ (điểm \%)} & \textbf{RPD (\%)} & \textbf{Đánh giá} \\
\hline
\endfirsthead
\multicolumn{7}{c}{\tablename\ \thetable{} -- tiếp theo} \\
\hline
\textbf{Ngày} & \textbf{KHM} & \textbf{Lần 1 (\(\%\mathrm{P_2O_5}\))} & \textbf{Lần 2 (\(\%\mathrm{P_2O_5}\))} & \textbf{$d$ (điểm \%)} & \textbf{RPD (\%)} & \textbf{Đánh giá} \\
\hline
\endhead
\hline
\multicolumn{7}{r}{\textit{Còn tiếp trang sau}} \\
\endfoot
\hline
\endlastfoot
01/06/2026 & 501X1 & 12,21 & 12,21 & 0,00 & 0,00 & Đạt \\
\hline
02/06/2026 & 03DGX1 & 16,52 & 16,38 & 0,14 & 0,85 & Đạt \\
\hline
03/06/2026 & 603X1 & 4,69 & 4,63 & 0,06 & 1,29 & Đạt \\
\hline
03/06/2026 & 401X2 & 8,41 & 8,53 & 0,12 & 1,42 & Đạt \\
\hline
03/06/2026 & 400X1 & 3,63 & 3,52 & 0,11 & 3,08 & Đạt \\
\hline
03/06/2026 & 300X1 & 5,91 & 6,05 & 0,14 & 2,34 & Đạt \\
\hline
03/06/2026 & 403X1 & 4,24 & 4,38 & 0,14 & 3,25 & Đạt \\
\hline
04/06/2026 & 401X1 & 9,40 & 9,40 & 0,00 & 0,00 & Đạt \\
\hline
04/06/2026 & 03DGX1 & 16,33 & 16,18 & 0,15 & 0,92 & Đạt \\
\hline
06/06/2026 & 12X1 & 5,16 & 4,92 & 0,24 & 4,76 & Đạt \\
\hline
\end{longtable}

}

\noindent\textbf{Nhận xét:} Trong tuần 01--06/06/2026, cả 10 cặp kết quả xác định lân hữu hiệu đều đạt yêu cầu RPD không quá 5\%.
Giá trị RPD dao động từ 0,00\% đến 4,76\%, với giá trị trung bình khoảng 1,79\%.
Mẫu 12X1 ngày 06/06/2026 có RPD lớn nhất (4,76\%) và nằm gần giới hạn chấp nhận, vì vậy cần tiếp tục theo dõi độ ổn định của các công đoạn tạo kết tủa, lọc, rửa, hòa tan kết tủa và nhận biết điểm kết thúc chuẩn độ trong các lần phân tích tiếp theo.

\subsection{Kết quả xác định carbon hữu cơ tổng số}
\CarbonResultsTable

\subsubsection{Đánh giá độ lặp lại}
Với phép Walkley--Black, chỉ tiêu được đánh giá là giá trị chất hữu cơ (OM). Bảng~\ref{tab:do-lap-lai-carbon} áp dụng tiêu chí nghiêm ngặt $\mathrm{RPD}\leq0{,}5\%$.

\begin{table}[H]
  \centering
  \small
  \renewcommand{\arraystretch}{1.25}
  \caption{Đánh giá độ lặp lại phép xác định chất hữu cơ theo từng mẫu}
  \label{tab:do-lap-lai-carbon}
  \resizebox{\textwidth}{!}{%
  \begin{tabular}{|>{\centering\arraybackslash}p{0.15\textwidth}|>{\centering\arraybackslash}p{0.14\textwidth}|>{\centering\arraybackslash}p{0.15\textwidth}|>{\centering\arraybackslash}p{0.15\textwidth}|>{\centering\arraybackslash}p{0.14\textwidth}|>{\centering\arraybackslash}p{0.14\textwidth}|>{\centering\arraybackslash}p{0.13\textwidth}|}
    \hline
    \textbf{Ngày} & \textbf{KHM} & \textbf{Lần 1 (\(\%\mathrm{OM}\))} & \textbf{Lần 2 (\(\%\mathrm{OM}\))} & \textbf{$d$ (điểm \%)} & \textbf{RPD (\%)} & \textbf{Đánh giá} \\
    \hline
    03/06/2026 & HC1X2 & 19,03 & 19,59 & 0,56 & 2,90 & Đạt \\
    \hline
    04/06/2026 & HC1X2 & 18,85 & 19,43 & 0,58 & 3,03 & Đạt \\
    \hline
    15/06/2026 & MHCX2 & 16,60 & 16,12 & 0,48 & 2,93 & Đạt \\
    \hline
    16/06/2026 & HC1X2 & 19,00 & 19,30 & 0,30 & 1,57 & Đạt \\
    \hline
    22/06/2026 & HC1X2 & 19,98 & 20,26 & 0,28 & 1,39 & Đạt \\
    \hline
  \end{tabular}
  }
\end{table}


\noindent\textbf{Nhận xét:} Cả 5 cặp kết quả OM đều có RPD lớn hơn 0,5\%, vì vậy chưa đạt yêu cầu độ lặp lại theo tiêu chí đã áp dụng. RPD dao động từ 1,92\% đến 12,33\%, trong đó chênh lệch lớn nhất thuộc mẫu HClX2 ngày 03/06/2026 có kết quả 21,91\% và 24,79\%. Các phép thử này cần được thực hiện lại sau khi tăng cường đồng nhất mẫu, kiểm tra thể tích dicromat, điều kiện oxy hóa và chuẩn độ bằng muối Mohr.

\subsection{Kết quả xác định độ ẩm}
\MoistureResultsTable

\subsubsection{Đánh giá độ lặp lại}
Bảng~\ref{tab:do-lap-lai-do-am} được lập theo cùng nguyên tắc để đánh giá độ lặp lại độ ẩm, với tiêu chí $\mathrm{RPD}\leq0{,}1\%$. Do chưa có số liệu cân và kết quả \%H$_2$O chính thức, các ô được giữ để bổ sung khi hoàn tất phép thử.

\begin{table}[H]
  \centering
  \small
  \renewcommand{\arraystretch}{1.3}
  \caption{Đánh giá độ lặp lại phép xác định độ ẩm theo từng mẫu}
  \label{tab:do-lap-lai-do-am}
  \resizebox{\textwidth}{!}{%
  \begin{tabular}{|>{\centering\arraybackslash}p{0.16\textwidth}|>{\centering\arraybackslash}p{0.14\textwidth}|>{\centering\arraybackslash}p{0.15\textwidth}|>{\centering\arraybackslash}p{0.15\textwidth}|>{\centering\arraybackslash}p{0.14\textwidth}|>{\centering\arraybackslash}p{0.13\textwidth}|>{\centering\arraybackslash}p{0.13\textwidth}|}
    \hline
    \textbf{Ngày} & \textbf{KHM} & \textbf{Lần 1 (\%H$_2$O)} & \textbf{Lần 2 (\%H$_2$O)} & \textbf{$d$ (điểm \%)} & \textbf{RPD (\%)} & \textbf{Đánh giá} \\
    \hline
    [Bổ sung] & [Bổ sung] & [Bổ sung] & [Bổ sung] & [Bổ sung] & [Bổ sung] & [Bổ sung] \\
    \hline
    [Bổ sung] & [Bổ sung] & [Bổ sung] & [Bổ sung] & [Bổ sung] & [Bổ sung] & [Bổ sung] \\
    \hline
  \end{tabular}
  }
\end{table}


\noindent\textbf{Nhận xét:} Sau khi bổ sung số liệu, cần tính RPD cho từng cặp kết quả \%H$_2$O. Mẫu chỉ được đánh giá đạt khi RPD không vượt quá 0,1\% và khối lượng sau các chu kỳ sấy đã đạt trạng thái không đổi. Chưa thể kết luận định lượng cho phép xác định độ ẩm ở thời điểm hiện tại.

\section{Đánh giá sai lệch giữa kết quả đo và giá trị công bố}
Giá trị phân tích trung bình của từng mẫu được so sánh với giá trị công bố trên nhãn, tiêu chuẩn cơ sở hoặc hồ sơ kỹ thuật tương ứng. Chênh lệch được tính như sau:
\[
  \Delta = \bar{x}-x_{\mathrm{cb}}, \qquad
  \delta(\%)=\frac{\Delta}{x_{\mathrm{cb}}}\times100.
\]
Trong đó, $\bar{x}$ là kết quả trung bình của các lần lặp hợp lệ, $x_{\mathrm{cb}}$ là giá trị công bố, $\Delta$ là sai lệch tuyệt đối và $\delta$ là sai lệch tương đối.

\begin{table}[H]
  \centering
  \small
  \renewcommand{\arraystretch}{1.3}
  \caption{Bảng đối chiếu kết quả đo với giá trị công bố}
  \label{tab:tong-hop-ket-qua-kiem-nghiem}
  \resizebox{\textwidth}{!}{%
  \begin{tabular}{|p{0.15\textwidth}|p{0.15\textwidth}|>{\centering\arraybackslash}p{0.13\textwidth}|>{\centering\arraybackslash}p{0.13\textwidth}|>{\centering\arraybackslash}p{0.12\textwidth}|>{\centering\arraybackslash}p{0.12\textwidth}|>{\centering\arraybackslash}p{0.12\textwidth}|}
    \hline
    \centering\textbf{Mã mẫu / sản phẩm} & \centering\textbf{Chỉ tiêu} & \textbf{Giá trị công bố (\%)} & \textbf{Kết quả TB (\%)} & \textbf{$\Delta$ (điểm \%)} & \textbf{$\delta$ (\%)} & \textbf{Nhận xét} \\
    \hline
    [Bổ sung] & Đạm tổng số & [Bổ sung] & [Bổ sung] & [Bổ sung] & [Bổ sung] & [Bổ sung] \\
    \hline
    [Bổ sung] & K$_2$O hữu hiệu & [Bổ sung] & [Bổ sung] & [Bổ sung] & [Bổ sung] & [Bổ sung] \\
    \hline
    [Bổ sung] & P$_2$O$_5$ hữu hiệu & [Bổ sung] & [Bổ sung] & [Bổ sung] & [Bổ sung] & [Bổ sung] \\
    \hline
    [Bổ sung] & Chất hữu cơ (OM) & [Bổ sung] & [Bổ sung] & [Bổ sung] & [Bổ sung] & [Bổ sung] \\
    \hline
    [Bổ sung] & Độ ẩm & [Bổ sung] & [Bổ sung] & [Bổ sung] & [Bổ sung] & [Bổ sung] \\
    \hline
  \end{tabular}
  }
\end{table}


\subsection{Nhận xét}
Các bảng số liệu hiện có cung cấp kết quả đo lặp, nhưng chưa kèm giá trị công bố chính thức của từng mã sản phẩm và chưa có số liệu độ ẩm. Vì vậy, chưa thể kết luận một mẫu đạt hay không đạt so với chỉ tiêu công bố. Sau khi bổ sung giá trị trên nhãn hoặc tiêu chuẩn kỹ thuật, cần tính $\Delta$, $\delta$ và đối chiếu với dung sai áp dụng cho đúng loại phân bón; chỉ nên đưa ra kết luận cuối cùng khi cả độ lặp lại và sai lệch với giá trị công bố đều đáp ứng yêu cầu.

\section{Nhận xét chung}
Các phép xác định đạm tổng số, kali hữu hiệu, lân hữu hiệu và chất hữu cơ đã có số liệu lặp để đánh giá bước đầu độ ổn định của thao tác. Kết quả kali cần luôn được gắn với đường chuẩn và mẫu trắng; kết quả lân và đạm cần kiểm soát điểm kết thúc chuẩn độ; còn kết quả carbon cần đặc biệt chú ý độ đồng nhất mẫu và hiệu quả oxy hóa. Số liệu độ ẩm và giá trị công bố của từng sản phẩm cần được bổ sung để hoàn tất đánh giá chất lượng toàn diện.

\section{Đề xuất xử lý khi kết quả chưa phù hợp}
Khi có chỉ tiêu không đạt hoặc kết quả lặp có độ phân tán lớn, cần xem xét theo Bảng~\ref{tab:de-xuat-xu-ly-ket-qua} trước khi đưa ra kết luận chính thức.

\begin{landscape}
\begingroup
\footnotesize
\setlength{\tabcolsep}{4pt}
\renewcommand{\arraystretch}{1.18}

\begin{longtable}{|>{\raggedright\arraybackslash}p{0.18\linewidth}|>{\raggedright\arraybackslash}p{0.20\linewidth}|>{\raggedright\arraybackslash}p{0.25\linewidth}|>{\raggedright\arraybackslash}p{0.29\linewidth}|}
\caption{Đề xuất xử lý các tình huống chưa phù hợp trong quá trình kiểm nghiệm}\label{tab:de-xuat-xu-ly-ket-qua} \\
\hline
\textbf{Tình huống} & \textbf{Biểu hiện} & \textbf{Nguyên nhân có thể} & \textbf{Đề xuất xử lý} \\
\hline
\endfirsthead

\multicolumn{4}{c}{\tablename~\thetable{} -- tiếp theo} \\
\hline
\textbf{Tình huống} & \textbf{Biểu hiện} & \textbf{Nguyên nhân có thể} & \textbf{Đề xuất xử lý} \\
\hline
\endhead

\hline
\multicolumn{4}{r}{\textit{Còn tiếp trang sau}} \\
\endfoot

\hline
\endlastfoot

Độ lặp không đạt yêu cầu
& Chênh lệch giữa hai lần phân tích lớn; RPD vượt giới hạn chấp nhận của phương pháp
& Mẫu chưa được nghiền và trộn đồng nhất trước khi chia mẫu phân tích
& Nghiền và đồng nhất lại mẫu; lấy lại hai phần mẫu độc lập để tiến hành phân tích song song. \\
\hline

&
& Sai số trong quá trình cân mẫu, hút pipet, định mức hoặc pha loãng
& Kiểm tra lại cân, pipet và dụng cụ định mức; thực hiện thao tác đúng quy trình và hạn chế thay đổi người thao tác giữa các lần lặp. \\
\hline

&
& Sai khác trong quá trình xử lý mẫu, thời gian phản ứng, nhiệt độ hoặc điều kiện thí nghiệm
& Kiểm soát thống nhất thời gian, nhiệt độ, lượng hóa chất và trình tự thao tác giữa các lần thử. \\
\hline

&
& Xác định điểm kết thúc chuẩn độ chưa đồng nhất
& Chuẩn độ chậm khi gần điểm cuối; sử dụng cùng loại chỉ thị và tiêu chí nhận biết màu; thực hiện mẫu trắng song song. \\
\hline

&
& Thiết bị phân tích hoạt động chưa ổn định, đặc biệt đối với máy quang kế ngọn lửa
& Kiểm tra đường chuẩn, dung dịch chuẩn kiểm soát, độ ổn định của ngọn lửa, hệ thống hút và đầu phun; vệ sinh thiết bị trước khi đo lại. \\
\hline

&
& Hóa chất hoặc dung dịch chuẩn có nồng độ không ổn định
& Kiểm tra hạn sử dụng, điều kiện bảo quản; chuẩn hóa lại dung dịch chuẩn khi cần thiết. \\
\hline

&
& --
& \textbf{Thực hiện lại toàn bộ phép thử từ bước cân mẫu. Chỉ lấy kết quả trung bình khi các lần phân tích mới đáp ứng yêu cầu độ lặp.} \\
\hline

Kết quả thành phẩm khác với chỉ tiêu công bố
& Hàm lượng N, \(\mathrm{P_2O_5}\), \(\mathrm{K_2O}\) hoặc chất hữu cơ thấp hơn giới hạn cho phép so với mức đăng ký
& Có thể xuất phát từ sai số phân tích hoặc quá trình chuẩn bị mẫu
& Trước tiên kiểm tra lại phép tính, hệ số pha loãng, đường chuẩn, mẫu trắng và điều kiện phân tích; sau đó lấy mẫu mới và phân tích lại để xác nhận. \\
\hline

&
& Thành phần dinh dưỡng thực tế của nguyên liệu khác với giá trị sử dụng khi tính công thức phối trộn
& Kiểm tra lại kết quả phân tích nguyên liệu đầu vào và hồ sơ chất lượng của từng lô nguyên liệu. \\
\hline

&
& Sai số trong quá trình định lượng nguyên liệu hoặc tỷ lệ phối trộn
& Kiểm tra phiếu phối liệu, hệ thống cân định lượng và lượng nguyên liệu thực tế đưa vào mẻ sản xuất. \\
\hline

&
& Hỗn hợp sau phối trộn chưa đồng đều
& Kiểm tra thời gian và điều kiện phối trộn; lấy mẫu tại nhiều vị trí của lô để đánh giá mức độ đồng nhất. \\
\hline

& Hàm lượng \(\mathrm{K_2O}\) hoặc \(\mathrm{P_2O_5}\) thấp bất thường
& Nguồn nguyên liệu kali hoặc lân chưa phù hợp, lượng nguyên liệu cấp vào dây chuyền sai lệch hoặc hỗn hợp phân bố không đồng đều
& Kiểm tra riêng nguồn nguyên liệu kali hoặc lân, lượng nguyên liệu được cấp vào dây chuyền và khả năng phân bố không đồng đều trong hỗn hợp. \\
\hline

& Độ ẩm thành phẩm cao hơn giới hạn
& Nguyên liệu có độ ẩm cao; sản phẩm hút ẩm trong quá trình sản xuất hoặc bảo quản; quá trình làm khô chưa phù hợp
& Kiểm tra độ ẩm nguyên liệu, điều kiện sấy hoặc làm khô, thời gian lưu sản phẩm và điều kiện kho; hạn chế tiếp xúc với môi trường có độ ẩm cao. \\
\hline

& Hàm lượng dinh dưỡng cao hơn mức đăng ký
& Có thể do sai số pha loãng, đường chuẩn hoặc lượng nguyên liệu dinh dưỡng thực tế cao hơn công thức
& Kiểm tra lại phép phân tích và công thức phối trộn; nếu kết quả phân tích lại vẫn tương tự thì đánh giá mức độ phù hợp theo giới hạn pháp lý áp dụng cho sản phẩm. \\
\hline

&
& --
& \textbf{Nếu kết quả phân tích lại vẫn không phù hợp, cần tạm thời chưa kết luận lô đạt chất lượng và chuyển thông tin cho Phòng Kỹ thuật/Bộ phận Sản xuất để xác định nguyên nhân, điều chỉnh quá trình sản xuất và kiểm nghiệm lại trước khi xuất sản phẩm.} \\

\end{longtable}
\endgroup
\end{landscape}

